\documentclass[10pt,a4paper]{amsart}
\usepackage[margin=22mm]{geometry}
\usepackage{fontspec,amsmath,amssymb,mathtools,mathalfa,tikz}
\setmainfont{Latin Modern Roman}[SmallCapsFont=lmromancaps10-regular.otf]
\usepackage[colorlinks=true,linkcolor=blue,citecolor=teal,urlcolor=blue]{hyperref}
\sloppy\newcommand{\eg}{e.g.\ }\newcommand{\ie}{i.e.\ }\newcommand{\resp}{resp.\ }
\usetikzlibrary{arrows,positioning,calc}
\numberwithin{equation}{section}
\renewcommand{\theenumi}{\roman{enumi}}
\hyphenation{dimen-sion depen-ding Essen-tial-ly}
\newcommand\sbw[2]{{{\mathsf{\Lambda}}^{{\!\scriptstyle#1}}{#2}}}

\DeclareMathOperator\codim{codim}
\DeclareMathOperator\Res{Res}

\newcommand\be{\boldsymbol{e}}
\newcommand\bx{\boldsymbol{x}}
\newcommand\by{\boldsymbol{y}}
\newcommand\bz{\boldsymbol{z}}
\newcommand\bX{\boldsymbol{X}}
\newcommand\bY{\boldsymbol{Y}}
\newcommand\bZZ{\boldsymbol{Z}}
\newcommand\bQ{\mathbb{Q}}
\newcommand\bN{\mathbb{N}}
\newcommand\bR{\mathbb{R}}
\newcommand\bZ{\mathbb{Z}}
\newcommand\cA{\mathcal{A}}
\newcommand\cB{\mathcal{B}}
\newcommand\cD{\mathfrak{D}}
\newcommand\cI{\mathcal{I}}
\newcommand\cU{\mathcal{U}}
\newcommand\cV{\mathcal{V}}
\newcommand\tP{\widetilde{P}}
\newcommand\tQ{\widetilde{Q}}
\newcommand\tR{\widetilde{R}}
\newcommand\tS{\widetilde{S}}
\newcommand\ee{\varepsilon}
\newcommand\GL{\mathrm{GL}}
\newcommand\GrO{\mathcal{O}}
\newcommand\hlambda{\widehat{\lambda}}
\newcommand\homega{\widehat{\omega}}\newcommand\HH{\mathcal{H}}
\newcommand\HHH{\mathcal{D}_\xi'}
\newcommand\HHstar{\mathcal{D}_{\xi}}
\newcommand{\norm}[1]{\left\|#1 \right\|_2}
\newcommand{\bnorm}[1]{\biggl\|#1 \biggr\|_2}
\newcommand{\normH}[1]{\|#1\|}
\newcommand{\psc}[2]{\left(#1 \mid #2\right)}
\newcommand{\Vect}[2][]{\left\la #2\right\ra_{#1}}

\let\subset \subseteq
\let\la\langle
\let\ra\rangle
\newcommand\AND{\quad\text{and}\quad}

\theoremstyle{definition}
\newtheorem{Def}{Definition}[section]
\newtheorem{Rem}[Def]{Remark}
\newtheorem*{remark*}{Remark}

\theoremstyle{plain}
\newtheorem{Prop}[Def]{Proposition}
\newtheorem{Lem}[Def]{Lemma}
\newtheorem{Thm}{Theorem}[section]
\newtheorem{Cor}[Def]{Corollary}
\begin{document}
\title{Uniform polynomial approximation in degrees at least three}
\author{Anthony Poëls}
\date{}
\maketitle
\section{Problem and original exponent notation}
Let $\xi$ be a non-zero real number and let $n$ be a positive integer. Dirichlet's theorem (1842) is one of the most basic results of Diophantine approximation. It shows that for any real number $H > 1$, there exists a non-zero integer point $(x_0,\dots,x_n)\in\bZ^{n+1}$ such that
\begin{equation}
\label{eq: Dirichlet system}
\max\big\{|x_1|,\dots,|x_n|\big\} \leq H \AND |x_0 + x_1\xi + \cdots + x_n\xi^n| \leq H^{-n}.
\end{equation}
It is natural to ask if we can improve the exponent $n$ of $H^{-n}$, and this question gives rise to two Diophantine exponents. The so-called \emph{uniform exponent of approximation} $\homega_n(\xi)$ (\resp the \emph{ordinary} exponent $\omega_n(\xi)$), is the supremum of the real numbers $\omega > 0$ such that the system
\begin{equation*}
\normH{P} \leq H \AND 0 <|P(\xi)| \leq H^{-\omega}
\end{equation*}
admits a non-zero solution $P\in\bZ[X]$ of degree at most $n$ for each sufficiently large $H$ (\resp for arbitrarily large $H$). Here, $\normH{P}$ denotes the (naive) \emph{height} of $P$, defined as the largest absolute value of its coefficients. These quantities have been extensively studied over the past half-century, see for example \cite{bugeaud2015exponents} for a nice overview of the subject. By Dirichlet's theorem, if $\xi$ is not an algebraic number of degree $\leq n$, then we have
\[
\omega_n(\xi) \geq \homega_n(\xi) \geq n,
\]
and it is well known that those inequalities are equalities for almost all real numbers~$\xi$ (with respect to the Lebesgue measure). Note that if $\xi$ is an algebraic number of degree $d$, then $\homega_n(\xi)$ and $\omega_n(\xi)$ are both equal to $\min\{n,d-1\}$ (it is a consequence of Schmidt's subspace theorem, see \cite[Th.\,2.10]{bugeaud2015exponents}). We can therefore restrict our study to the set of transcendental real numbers. The initial question ``can we improve the exponent $n$ in Dirichlet's theorem?'' may be rephrased as follows: ``does there exist a transcendental real number $\xi$ satisfying $\homega_n(\xi) > n$?''. For $n=1$ the answer is negative and rather elementary to prove, so the first non-trivial case is $n=2$. Before the early $2000$s, it was conjectured that no such number existed. This belief was swept away by Roy's extremal numbers \cite{roy2003approximation}, \cite{roy2004approximation}, \cite{arbourRoyCriterionDegreTwo}, whose exponent $\homega_2$ is equal to the maximal possible value $(3+\sqrt 5)/2 = 2.618\cdots$. Since then, several families of transcendental real numbers whose uniform exponent $\homega_2$ is greater than $2$ have been discovered (see for example \cite{roy2007two}, \cite{bugeaud2005exponentsSturmian}, \cite{poels2017exponents, poelsExpoGeneralClass2021}). However, for $n \geq 3$ the mystery remains, and it is still an open question whether or not there exists $\xi\in\bR\setminus\overline{\bQ}$ with $\homega_n(\xi) > n$.
\begin{Thm}
\label{Thm : main d=2}
Let $n\geq 3$ be an integer and $\xi\in\bR$ be a transcendental real number. If $n\geq 4$, then
\begin{equation*}
\homega_n(\xi) \leq 2n-2.
\end{equation*}
For $n=3$, we have the weaker estimate $\homega_3(\xi)\leq 2+\sqrt 5 = 4.23\cdots$.
\end{Thm}
\setcounter{section}{1}
\section{Notation}
\label{section: notation}
Throughout this paper, $\xi$ denotes a transcendental real number. The floor (\resp ceiling) function is denoted by $\lfloor \cdot \rfloor$ (\resp $\lceil \cdot \rceil$). If $f, g : I \to [0, +\infty)$ are two functions on a set $I$, we write $f = \GrO(g)$ or $f \ll g$ or $g \gg f$ to mean that there is a positive constant $c$ such that $f (x) \leq cg(x)$ for each $x\in I$. We write $f \asymp g$ when both $f \ll g$ and $g \ll f$ hold.

Let $K$ be a field. If $\cA$ is a subset of a $K$-vector space $V$, we denote by $\Vect[K]{\cA} \subset V$ the $K$-vector space spanned by $\cA$, with the convention that $\Vect[K]{\emptyset} = \{0\}$.

Given a ring $A$ (typically $A=\bR$ or $\bZ$) and an integer $n\geq 0$, we denote by $A[X]$ the ring of polynomials in $X$ with coefficients in $A$, and by $A[X]_{\leq n} \subset A[X]$ the subgroup of polynomials of degree at most $n$. We say that $P\in\bZ[X]$ is \emph{primitive} if it non-zero and the greatest common divisor of its coefficients is $1$. Given $P=\sum_{k=0}^{n} a_kX^k \in \bR[X]$, we set
\begin{equation*}
\normH{P} = \max_{0\leq k \leq n} |a_k|.
\end{equation*}

Gelfond's lemma is the following statement (see \eg \cite[Lem.\,A.3]{bugeaud2004approximation} as well as \cite{brownawell1974sequences}). For any non-zero polynomials $P_1,\dots,P_r\in\bR[X]$ with product $P = P_1\cdots P_r$ of degree at most $n$, we have
\begin{equation}
\label{eq: Gelfond's lemma}
e^{-n} \normH{P_1}\cdots \normH{P_r} < \normH{P} < e^{n} \normH{P_1}\cdots \normH{P_r}.
\end{equation}
In particular, for each non-zero polynomial $P \in \bZ[X]_{\leq n}$ and each factor $Q\in\bZ[X]$ of $P$, we have $e^{-n}\normH{Q} < \normH{P}$. We will often use \eqref{eq: Gelfond's lemma} as follows. If $Q\in\bZ[X]_{\leq n}$ is irreducible and if $P\in\bZ[X]_{\leq n}$ is a non-zero polynomial which satisfies $\normH{P} \leq e^{-n}\normH{Q}$, then $Q$ cannot divide $P$. They are thus coprime polynomials.

We recall the definition of the resultant, which, as explained in the introduction, is useful for estimating the exponent $\homega_n(\xi)$ (see also Section~\ref{section: resultant et premier resultat}). Let $P,Q\in\bZ[X]$ be non-constant polynomials of degree $p$ and $q$ respectively, and let $a_i,b_j\in \bZ$ such that $P(X) = \sum_{k=0}^{p} a_kX^k$ and $Q(X) = \sum_{k=0}^{q} b_kX^k$.
Their \emph{resultant} $\Res(P,Q)$ is defined as the $(q+p)$-dimensional determinant
\begin{equation}
\label{eq det resultant}
\arraycolsep3.5pt
\Res(P,Q) = \begin{array}{cc}
\left | \begin{array}{cccc}
a_p & 0 & \dots \\
a_{p-1} & a_p & \\
\vdots & \vdots & \ddots \\
a_0 & \\
0 & a_0 \\
\vdots & \vdots & \ddots \\
& & & a_0
\end{array}\right. &
\left. \hspace*{-10pt}
\begin{array}{cccc}
b_q & 0 & \dots \\
b_{q-1} & b_q & \\
\vdots & \vdots & \ddots \\
b_0 & \\
0 & b_0 \\
\vdots & \vdots & \ddots \\
& & & b_p
\end{array}\right| \\
\underbrace{\hspace*{2.7cm}}_{q} &
\underbrace{\hspace*{2.3cm}}_{p}
\end{array}.
\end{equation}
Besides the exponents of linear approximation $\omega_n$ and $\homega_n$, we will also need the following exponents of simultaneous rational approximation. For each positive integer~$n$, the exponent $\hlambda_n(\xi)$ (\resp $\lambda_n(\xi)$) is the supremum of the real numbers $\lambda \geq 0$ such that the system
\begin{equation*}
|y_0| \leq Y \AND L(\by) \leq Y^{-\lambda} \quad \textrm{where } L(\by):=\max_{1\le k \le n} |y_0\xi^k - y_k|,
\end{equation*}
admits a non-zero integer solution $\by=(y_0,\dots,y_n)\in\bZ^{n+1}$ for each sufficiently large $Y\geq 1$ (\resp for arbitrarily large $Y$). Dirichlet's theorem \cite[\S II.1, Th.\,1A]{schmidt1996diophantine} implies that $\hlambda_n(\xi)\geq 1/n$. The best upper bounds known to date for $\hlambda_n(\xi)$ when $n\geq 4$ are established in joint work with Roy in \cite{poels2022simultaneous}. In particular, there is an explicit positive constant $a$ such that
\[
\hlambda_n(\xi) \leq \frac{1}{n/2+an^{1/2}+1/3},
\]
and sharper results are also obtained when $n$ is small.

\section{Minimal polynomials}
\label{section: minimal pol}

A \emph{sequence of minimal polynomials} (associated to $n$ and $\xi$) is a sequence $(P_i)_{i\geq 0}$ of non-zero polynomials in $\bZ[X]_{\leq n}$ satisfying the following properties:
\begin{enumerate}
\item the sequence $\big(\normH{P_i}\big)_{i\geq 0}$ is strictly increasing;
\item the sequence $\big(|P_i(\xi)|\big)_{i\geq 0}$ is strictly decreasing;
\item if $|P(\xi)| < |P_i(\xi)|$ for some index $i\geq 0$ and a non-zero $P\in\bZ[X]_{\leq n}$, then $\normH{P}\geq \normH{P_{i+1}}$.
\end{enumerate}

Note that if we require the dominant coefficient of $P_i$ to be positive (and since $\xi$ is transcendental), then the above sequence is unique up to its first terms. Let $(P_i)_{i\geq 0}$ be a sequence as above. We have the classical formulas:
\begin{equation}
\label{eq: exposant via minimal points}
\homega_n(\xi) = \liminf_{i\to \infty} \frac{-\log |P_i(\xi)|}{\log \normH{P_{i+1}}} \AND \omega_n(\xi) = \limsup_{i\to \infty} \frac{-\log |P_i(\xi)|}{\log \normH{P_{i}}}.
\end{equation}
In particular, given a positive real number $\homega$ with $\homega < \homega_n(\xi)$, then we have, for each sufficiently large index $i$,
\begin{equation}
\label{eq: H(P_i+1) controle par H(P_i)}
|P_i(\xi)| \leq \normH{P_{i+1}}^{-\homega} \AND \normH{P_{i+1}}^\tau \leq \normH{P_i},\quad \textrm{where } \tau:= \frac{\homega}{\omega_n(\xi)},
\end{equation}
(with the convention $\tau = 0$ if $\omega_n(\xi) = \infty$). The second inequality in \eqref{eq: H(P_i+1) controle par H(P_i)} asks for an upper bound on $\omega_n(\xi)$. Given a non-zero $P \in \bZ[X]$, we set $\omega(P)=0$ if $\normH{P} = 1$. Otherwise, we denote by $\omega(P)$ the real number satisfying
\[
|P(\xi)| = \normH{P}^{-\omega(P)}.
\]
With this notation, we have
\begin{equation}
\label{eq: encadrement omega(P_i)}
\omega_n(\xi) = \limsup_{\substack{\normH{P}\to \infty \\ P\in\bZ[X]_{\leq n}}} \omega(P) = \limsup_{i\to \infty} \omega(P_i) \AND \liminf_{i\to \infty} \omega(P_i) \geq \homega_n(\xi).
\end{equation}
The following results are well-known. We prove them for the sake of completion. The first one follows from the arguments of the proof of \cite[Lem.\,2]{davenport1967approximation} (see also \cite[Lem.\,4.1]{roy2004approximation}).

\begin{Lem}
\label{Lem: P_i, P_i+1 est une Z-base}
Let $i\geq 0$ and write $V_i = \Vect[\bR]{P_i,P_{i+1}} \subset \bR[X]_{\leq n}$. Then $\{P_i, P_{i+1}\}$ forms a $\bZ$-basis of the lattice $V_i\cap \bZ[X]_{\leq n}$.
\end{Lem}

\begin{proof}
By contradiction, suppose that $\{P_i, P_{i+1}\}$ is not a $\bZ$-basis of $V_i\cap \bZ[X]_{\leq n}$. Then there exists a non-zero $Q\in \bZ[X]_{\leq n}$ which may be written as $Q = rP_i + sP_{i+1}$, where $r,s\in\bQ$ satisfy $|r|,|s|\leq 1/2$. In particular, we have
\begin{align*}
\normH{Q} &\leq |r|\normH{P_i} + |s|\normH{P_{i+1}} < \normH{P_{i+1}}, \\
|Q(\xi)| &\leq |r||P_i(\xi)| + |s||P_{i+1}(\xi)| < |P_{i}(\xi)|.
\end{align*}
This contradicts the minimality property of $P_i$.
\end{proof}

The next result is analogous to the second part of \cite[Lem.\,4.1]{roy2004approximation}. The construction of $S_i$ is due to Davenport and Schmidt \cite{davenport1967approximation}.

\begin{Lem}
\label{Lem: P_iH_i+1 = P_j-1H_j}
For each $i\geq 0$, define
\[
S_i = P_i(\xi)P_{i+1} - P_{i+1}(\xi)P_i \in\bR[X]_{\leq n}.
\]
Then
\[
\frac{1}{2}\normH{S_i} \leq \normH{P_{i+1}}|P_i(\xi)| \leq 2\normH{S_i}.
\]
Moreover, if for integers $0\leq i < j$ the space spanned by $P_i, P_{i+1}, \cdots, P_{j}$ has dimension~$2$, then $S_{j-1} = \pm S_i$. In particular
\[
\normH{P_{i+1}}|P_i(\xi)| \asymp \normH{P_{j}}|P_{j-1}(\xi)|.
\]
\end{Lem}

\begin{Rem}
Note that the quantity $\normH{S_i}$ satisfies $\normH{S_i} \asymp \HHstar(V_i)$, where $\HHstar$ is defined in Section~\href{https://jep.centre-mersenne.org/item/10.5802/jep.265.pdf#page=20}{8.2} and $V_i=\Vect[\bR]{P_i,P_{i+1}}$. We will study the function $\HHstar$ more deeply later.
\end{Rem}

\begin{proof}
We easily get $\normH{S_i} \leq 2\normH{P_{i+1}}|P_i(\xi)|$. Define $R_+, R_- \in\bZ[X]_{\leq n}$ by
\[
R_\pm = P_{i+1} \pm P_i.
\]
Suppose that there exists $\ee\in\{+,-\}$ such that $|R_\ee(\xi)| \leq |P_i(\xi)|/2$. Then by minimality of $P_i$, we must have $\normH{R_\ee} \geq \normH{P_{i+1}}$. Since $S_i = P_i(\xi)R_\ee - R_\ee(\xi) P_i$, we find
\begin{equation*}
\normH{S_i} \geq |P_i(\xi)|\normH{R_\ee} - |R_\ee(\xi)|\normH{P_i} \geq \frac{1}{2} \normH{P_{i+1}}|P_i(\xi)|.
\end{equation*}
Assume that $|R_+(\xi)|, |R_-(\xi)| \geq |P_i(\xi)|/2$. This is equivalent to
\[
|P_{i+1}(\xi)| \leq \frac{1}{2}|P_i(\xi)|.
\]
Again, this yields $\normH{S_i} \geq |P_i(\xi)|\normH{P_{i+1}} - |P_{i+1}(\xi)|\normH{P_i} \geq \normH{P_{i+1}}|P_i(\xi)|/2$.

Now, let us write $V_i = \Vect[\bR]{P_i,\dots,P_j}$, with $j > i$, and suppose that $V_i$ has dimension~$2$. We need to prove that $S_{j-1} = \pm S_i$. If $j=i+1$ it is automatic, we may therefore assume that $j\geq i+2$. By Lemma~\ref{Lem: P_i, P_i+1 est une Z-base}, there exist $a,b\in\bZ$ such that $P_i = aP_{i+1} + bP_{i+2}$. Since $\{P_i,P_{i+1}\}$ is also a $\bZ$-basis of $V_i$, we have $b=\pm 1$, and we deduce that
\[
S_i = \big(aP_{i+1}(\xi) + bP_{i+2}(\xi)\big)P_{i+1} - P_{i+1}(\xi)\big(aP_{i+1} + bP_{i+2}\big) = -bS_{i+1} = \pm S_{i+1}.
\]
By induction, we get $S_i = \pm S_{i+1} = \cdots = \pm S_{j-1}$.
\end{proof}

The proof of \cite[Lem.\,3]{davenport1967approximation} (which deals with the case $n=2$) yields the classical following result.

\begin{Lem}
\label{lem: I infini}
Suppose $n\geq 2$. Then, there are infinitely many indices $i\geq 1$ for which $P_{i-1}$, $P_i$ and $P_{i+1}$ are linearly independent.
\end{Lem}

\begin{proof}
By contradiction, suppose that there exists $i\geq 0$ such that $V=\Vect[\bR]{P_i,P_{i+1},\dots}$ has dimension $2$. By Lemma~\ref{Lem: P_iH_i+1 = P_j-1H_j} there exists $c>0$ such that for each $j > i$ we have
\[
0 < \normH{P_{i+1}}|P_i(\xi)| \leq c\normH{P_j}|P_{j-1}|.
\]
This leads to a contradiction since $\normH{P_j}|P_{j-1}| \leq \normH{P_j}^{1-\homega_n(\xi)+o(1)}$ tends to $0$ as $j$ tends to infinity.
\end{proof}

\begin{Rem}
As mentioned in the introduction, it is however possible that all polynomials $P_i$ with $i$ large enough lie in a subspace of dimension $3$ , see \cite[Th.\,1.3]{moshchevitin2007best}.
\end{Rem}

\section{Resultant and first estimates}
\label{section: resultant et premier resultat}

The following useful result can easily be derived from the proof of \cite[\S 5]{davenport1969approximation} (see also of \cite[Lem.\,1]{brownawell1974sequences}). We recall the arguments since they illustrate (in a simpler situation) how we will deal with generalized determinants.

\begin{Lem}
\label{lem: estimation classique du resultant}
Let $p,q$ be positive integers with $p,q\leq n$. There exists a constant $c>0$ depending on $\xi$ and $n$ only, with the following property. For any polynomials $P,Q\in\bZ[X]$ of degree $p$ and $q$ respectively, we have
\begin{equation*}
|\Res(P,Q)| \leq c \normH{P}^{q-1}\normH{Q}^{p-1} \max\big\{\normH{P}|Q(\xi)|, \normH{Q}|P(\xi)| \big\}.
\end{equation*}
\end{Lem}

\begin{proof}
Let $a_i,b_j\in \bZ$ such that $P(X) = \sum_{k=0}^{p} a_kX^k$ and $Q(X) = \sum_{k=0}^{q} b_kX^k$. For $i=1,\dots,p+q-1$, we add to the last row of the determinant \eqref{eq det resultant} its $i$-th row multiplied by $\xi^{p+q-i}$. This last row now becomes
\begin{equation*}
\Big(\xi^{q-1}P(\xi), \dots, \xi P(\xi), P(\xi), \xi^{p-1}Q(\xi), \dots, \xi Q(\xi), Q(\xi)\Big).
\end{equation*}
Using the upper bounds $|a_i| \leq \normH{P}$ and $|b_j|\leq \normH{Q}$ for the other entries of \eqref{eq det resultant}, we~obtain
\begin{equation*}
|\Res(P,Q)| \ll \normH{P}^{q-1}|P(\xi)| \normH{Q}^p + \normH{P}^q\normH{Q}^{p-1}|Q(\xi)|,
\end{equation*}
where the implicit constant only depends on $p, q$ and $\xi$.
\end{proof}

The next result, which is also based on inequalities involving resultants, will be used in Section~\href{https://jep.centre-mersenne.org/item/10.5802/jep.265.pdf#page=27}{10}. It ensures that if $R\in\bZ[X]$ is a ``good'' approximation, in the sense that $R(\xi)$ is very small compared to $\normH{R}$, and if we write $R$ as a product of coprime polynomials $B_1\cdots B_k$, then one of those factors is also a ``good'' approximation, while the product of the others is not.

\begin{Lem}
\label{lem: lem general 2}
Let $m,k$ be positive integers. There exists a constant $c>0$ depending on $m$ and $\xi$ only, with the following property. Let $B_1,\dots,B_k\in\bZ[X]$ be non-constant, pairwise coprime polynomials, and suppose that $R:= B_1\cdots B_k$ has degree at most $m$. Then, there exists $j\in\{1,\dots,k\}$ such that
\[
|B_j(\xi)| \leq c\normH{R}^{m-1}|R(\xi)| \AND \prod_{\substack{i=1 \\ i\neq j}}^{k} |B_i(\xi)| \geq c^{-1}\normH{R}^{-(m-1)}.
\]
\end{Lem}

\begin{proof}
If $k=1$ this is trivial. We now suppose that $k\geq 2$ and write $d_i = \deg(B_i)$ for $i=1,\dots, k$. By hypothesis, we have $\deg(R) = d_1+\cdots + d_k \leq m$. Note that
\begin{equation}
\label{eq proof: log H(P) = sum log H(R_i)}
|R(\xi)| = \prod_{i=1}^{k} |B_i(\xi)| \AND \normH{R} \asymp \prod_{i=1}^{k} \normH{B_i},
\end{equation}
the second inequality coming from Gelfond's lemma (with an implicit constant depending only on $m$). Choose $j\in\{1,\dots,k\}$ such that $|B_j(\xi)|$ is minimal and fix $i\in\{1,\dots,k\}$ with $i\neq j$. Since $B_i$ and $B_j$ are coprime, their resultant $\Res(B_i,B_j)$ is a non-zero integer. Using Lemma~\ref{lem: estimation classique du resultant}, we find
\begin{align*}
1 \leq |\Res(B_i,B_j)| & \ll \normH{B_i}^{d_j-1} \normH{B_j}^{d_i-1}\big( \normH{B_j}|B_i(\xi)| + \normH{B_i}|B_j(\xi)|\big) \\
& \ll \normH{B_i}^{d_j} \normH{B_j}^{d_i}|B_i(\xi)|,
\end{align*}
with an implicit constant depending only on $\xi$ and $m$, hence
\[
-\log |B_i(\xi)| \leq d_j\log \normH{B_i} + d_i\log \normH{B_j} + \GrO(1).
\]
On the other hand, by summing the above inequalities for $i\neq j$, and by using \eqref{eq proof: log H(P) = sum log H(R_i)}, we obtain
\begin{align*}
\sum_{\substack{i=1 \\ i\neq j}}^{k} -\log |B_i(\xi)|
& \leq d_j\sum_{\substack{i=1 \\ i\neq j}}^{k} \log \normH{B_i} +(m-d_j) \log \normH{B_j} + \GrO(1) \\
& \leq (m-1)\log \normH{R} + \GrO(1).
\end{align*}
We easily deduce that
\[
\prod_{\substack{i=1 \\ i\neq j}}^{k} |B_i(\xi)| \gg \normH{R}^{-(m-1)} \AND
|R(\xi)| = \prod_{i=1}^{k}|B_i(\xi)| \gg |B_j(\xi)| \normH{R}^{- (m-1)}.\qedhere
\]
\end{proof}

\section{A sequence of irreducible polynomials}
\label{section: sequence des Q_i}

As explained in the introduction, to get the upper bound $\homega_n(\xi) \leq 2n-1$, the strategy of Davenport and Schmidt \cite{davenport1969approximation} consists in considering the resultant $\Res(P,Q)$ of two ``good'' polynomial approximations $P,Q\in \bZ[X]_{\leq n}$. To ensure that $\Res(P,Q)$ does not vanish, they need a polynomial $P$ which is irreducible (for it is then easy to find $Q$ so that $P$ and $Q$ are coprime). The same difficulty appears in \cite{bugeaudSchleischitz2016Uniform}. Similarly, we will not work directly with a sequence of minimal polynomials. Instead, we will considerer the largest irreducible factors of the minimal polynomials. Now, let $n,d$ be integers with
\[
2\leq d < 1+ \frac{n}{2}.
\]
In this section, we assume that the transcendental real number $\xi$ satisfies $\homega_n(\xi) > 2n -d$ and we fix a real number $\homega$ (arbitrarily close to $\homega_n(\xi)$) such that
\begin{equation}
\label{eq: notation homega}
\homega_n(\xi) > \homega > 2n -d.
\end{equation}

We denote by $(P_i)_{i\geq 0}$ a sequence of minimal polynomials associated to $n$ and $\xi$. Our goal is to prove the existence of a sequence $(Q_i)_{i\geq 0}$ as below.

\begin{Prop}
\label{prop: existence Qi}
Suppose that \eqref{eq: notation homega} holds. Then, there exist a sequence $(Q_i)_{i\geq 0}$ of pairwise distinct polynomials in $\bZ[X]_{\leq n}$ and an index $j_0\geq 0$ with the following properties. The sequence $(\normH{Q_i})_{i\geq 0}$ is bounded below by $2$, unbounded and non-decreasing, and for any $i\geq 0$
\begin{enumerate}
\item \label{enum: property (Q_i) item 1} $Q_i$ is irreducible (over $\bZ$) and has degree at least $n-d+2$;
\item $Q_i$ divides $P_j$ for some index $j\geq j_0$ (not necessarily unique), and for each $j\geq j_0$ there exists $k\geq 0$ such that $Q_k$ divides $P_j$;
\item \label{enum: property (Q_i) item 6} $|Q_i(\xi)| = \normH{Q_i}^{-\omega(Q_i)}\leq \normH{Q_i}^{-\homega}$, and we further have
\begin{equation}
\label{eq: omega_n donné par les Q_i}
\omega_n(\xi) = \limsup_{k\to \infty} \omega(Q_k) \AND \liminf_{k\to \infty} \omega(Q_k) \geq \homega_n(\xi).
\end{equation}
\item \label{enum: H(P_j) controle par H(Q_i)} if $Q_i$ divides a minimal polynomial $P_j$ with $j\geq j_0$, then
\begin{equation}
\label{eq: premiere estimation Pj vs Qi}
\normH{P_j} \leq \normH{Q_i}^{1+\theta_i}, \quad \textrm{where } \theta_i = \frac{\omega(Q_i)-2n+d}{n-2d+3};
\end{equation}
\item \label{enum: H(Q_i+1) controle par H(Q_i)} we have
\begin{equation}
\label{eq: estimation Q_i+1 par Q_i}
\normH{Q_{i+1}}^\tau \leq \normH{Q_i} \quad \textrm{where } \tau = \frac{\homega\big(\homega -n-d+3\big)}{\omega_n(\xi)\big(\omega_n(\xi)-n-d+3\big)},
\end{equation}
with the convention $\tau= 0$ if $\omega_n(\xi) = \infty$.
\end{enumerate}
\end{Prop}

The above proposition is essentially a consequence of Lemma~\ref{lem: existence facteur irreductible grand} below. Assertion~\eqref{enum: property (Q_i) item 6} ensures that the polynomials $Q_i$ are quite good approximations, and they can be used to compute the exponent of best approximation $\omega_n(\xi)$. Estimate \eqref{eq: estimation Q_i+1 par Q_i} is the analog of the second inequality of \eqref{eq: H(P_i+1) controle par H(P_i)} but is way more difficult to prove. The main reason behind this difficulty is that there may be many polynomials $P\in\bZ[X]_{\leq n}$ with $\normH{Q_i} < \normH{P} < \normH{Q_{i+1}}$ and $|P(\xi)| < Q_i(\xi)$.

In order to prove Proposition~\ref{prop: existence Qi}, we need the two technical lemmas below. Essentially, they will be used to prove that the factors of $P_i$ of small degree are bad approximations. This will lead to the existence of a factor of large degree which is necessarily a rather good approximation.

\begin{Lem}
\label{lemme: omega(R) petit degre}
Suppose that \eqref{eq: notation homega} holds. Then, there exists a constant $c\in(0,1)$ depending only on $\xi$ and $n$ such that for any non-zero polynomial $R\in \bZ[X]_{\leq n-d+1}$ we have
\begin{equation}
\label{lemme; eq1 : omega(R) petit degree}
|R(\xi)| \geq c\normH{R}^{-(n + \deg(R) - 1)} \geq c\normH{R}^{-(2n-d)}.
\end{equation}
In particular \eqref{lemme; eq1 : omega(R) petit degree} holds for any $R\in\bZ[X]_{\leq d-2}$.
\end{Lem}

\begin{proof}
If $R$ is constant we have $|R(\xi)| = \normH{R}$ and the result is trivial. Now, suppose that $R$ is irreducible and not constant. We adapt the arguments of Davenport and Schmidt \cite[\S 5--6]{davenport1969approximation}. Set $H = e^{-n}\normH{R}$. By definition of $\homega_n(\xi)$ and $\homega$, if $H$ is sufficiently large, there exists a non-zero $P\in\bZ[X]_{\leq n}$ such that
\[
\normH{P} \leq H \AND |P(\xi)| \leq H^{-\homega}.
\]
By \eqref{eq: Gelfond's lemma}, the (irreducible) polynomial $R$ is not a factor of $P$, they are thus coprime polynomials. Their resultant is a non-zero integer, and using Lemma~\ref{lem: estimation classique du resultant}, we obtain
\begin{align*}
1 & \ll \normH{P}^{\deg(R)-1}\normH{R}^{n}|P(\xi)| + \normH{P}^{\deg(R)}\normH{R}^{n-1}|R(\xi)| \\
& \ll H^{n+\deg(R)-1-\homega} + H^{n+\deg(R)-1}|R(\xi)|.
\end{align*}
Since $\homega > 2n-d$ and $\deg(R) \leq n-d+1$, the first term tends to $0$ as $H$ tends to infinity. Hence $ 1 \ll H^{n+\deg(R)-1}|R(\xi)|$, which implies \eqref{lemme; eq1 : omega(R) petit degree}.

If $R$ is not irreducible, we write $R = \prod_{i=1}^{s}R_i$ with integer $s\geq 1$ and $R_1,\dots,R_s\in\bZ[X]$ irreducible of degree $\leq \deg(R)$ (possibly constant). Combining $\normH{R} \asymp \prod_{i=1}^{s}\normH{R_i}$ together with \eqref{lemme; eq1 : omega(R) petit degree} applied to the irreducible factors $R_i$, we find
\begin{equation*}
|R(\xi)| = \prod_{i=1}^{s}|R_i(\xi)| \gg \prod_{i=1}^{s}\normH{R_i}^{-(n + \deg(R) - 1)} \gg \normH{R}^{-(n + \deg(R) - 1)}.
\end{equation*}
Finally, the last assertion comes from the fact that $d-1\leq n+d-1$ (since $d\leq 1+n/2$).
\end{proof}

\begin{Lem}
\label{lem: existence facteur irreductible grand}
Suppose that \eqref{eq: notation homega} holds. There exist $i_0\geq 0$ and a constant $c>0$ such that for each $i\geq i_0$ the polynomial $P_i$ has a unique irreducible factor $\tP_i\in\bZ[X]$ of degree $\geq n-d+2$ and positive leading coefficient. It satisfies
\begin{equation}
\label{eq lem: existence facteur irr grand eq 2}
|P_i(\xi)| \normH{P_i}^{n+d-3} \geq c|\tP_i(\xi)| \normH{\tP_i}^{n+d-3},
\end{equation}
moreover $\big(\normH{\tP_i}\big)_{i\geq i_0}$ tends to infinity and as $i$ tends to infinity. For each $i$ large enough, we have $\normH{\tP_i}>1$, and writing $|\tP_i(\xi)| = \normH{\tP_i}^{-\omega(\tP_i)}$, we furthermore have
\begin{equation}
\label{eq lem: existence facteur irr grand eq 1}
\omega_n(\xi) = \limsup_{i\to\infty} \omega(\tP_i) \AND \liminf_{i\to\infty} \omega(\tP_i) \geq \homega_n(\xi).
\end{equation}
\end{Lem}

\begin{proof}
First, note that since $d < 1+n/2$, if we decompose $P_i$ as a product of irreducibles, there is at most one factor of degree $\geq n-d+2$. Fix $i\geq 0$ large enough so that $\omega(P_i) \geq \homega$, and write
\[
P:=P_i = \prod_{k=1}^s R_k,
\]
where $R_1,\dots,R_s\in\bZ[X]$ are irreducible polynomials (and $s$ is a positive integer). Suppose that $\deg(R_k) \leq n-d+1$ for each $k=1,\dots,s$. Then, by Lemma~\ref{lemme: omega(R) petit degre} together with $\normH{P} \asymp \prod_k \normH{R_k}$, we find
\begin{equation*}
\normH{P}^{-\homega}\geq |P(\xi)| = \prod_{k=1}^{s} |R_k(\xi)| \gg \prod_{k=1}^{s} \normH{R_k}^{-(2n-d)} \asymp \normH{P}^{-(2n-d)}.
\end{equation*}
This is impossible if $i$ is sufficiently large since $\homega > 2n-d$. Therefore, if $i$ is large enough, one of the factors $R_k$ has degree at least $n-d+2$. Without loss of generality, we may suppose that it is $R:=R_1$.
Write $S:= \prod_{k = 2}^{s} R_k$, so that $P=RS$. We have $\deg(S) \leq d-2$, and \eqref{lemme; eq1 : omega(R) petit degree} of Lemma~\ref{lemme: omega(R) petit degre} yields
\begin{equation*}
|S(\xi)| \gg \normH{S}^{-(n+d-3)}.
\end{equation*}
Together with $\normH{P}\asymp \normH{R}\normH{S}$, this leads to
\begin{equation*}
|P(\xi)| = |R(\xi)||S(\xi)| \gg |R(\xi)|\normH{S}^{-(n+d-3)} \asymp |R(\xi)|\normH{R}^{n+d-3}\normH{P}^{-(n+d-3)},
\end{equation*}
and \eqref{eq lem: existence facteur irr grand eq 2} follows easily by setting $\tP_i:=R$. The rest of the proof is based only on \eqref{eq lem: existence facteur irr grand eq 2} and the inequality $\normH{\tP_i} \ll \normH{P_i}$. Note that $|P_i(\xi)|\normH{P_i}^{n+d-3} \ll \normH{P_i}^{n+d-3-\homega}$ tends to $0$ as $i$ tends to infinity (using $d < 1+n/2$ together with $\homega > 2n-d$). We deduce that $|\tP_i(\xi)|\normH{\tP_i}^{n+d-3}$ also tends to $0$ as $i$ tends to infinity, which is possible only if $\normH{\tP_i}$ tends to infinity. In particular, if $i$ is large enough we must have $\normH{\tP_i} > 1$. Writing $|\tP_i(\xi)| = \normH{\tP_i}^{-\omega(\tP_i)}$, we also have $\omega(\tP_i) > n+d-3$. Now, using $\normH{\tP_i}\ll \normH{P_i}$, and taking the logarithms of the two sides of \eqref{eq lem: existence facteur irr grand eq 2}, we get
\begin{align*}
\big(\omega(P_i)-(n+d-3)\big)\log \normH{P_i} & \leq \big(\omega(\tP_i)-(n+d-3)\big)\log \normH{\tP_i} + \GrO(1) \\
& \leq \big(\omega(\tP_i)-(n+d-3)\big)\big(\log \normH{P_i}+\GrO(1)\big) + \GrO(1).
\end{align*}
By dividing by $\log \normH{P_i}$ and by simplifying, we deduce that $\omega(\tP_i) \geq \omega(P_i)\big (1 - o(1)\big)$ and \eqref{eq lem: existence facteur irr grand eq 1} follows easily from \eqref{eq: encadrement omega(P_i)}.
\end{proof}

\begin{proof}[Proof of Proposition~\ref{prop: existence Qi}]
Let $i_0\geq 0$ and $(\tP_i)_{i\geq i_0}$ given by Lemma~\ref{lem: existence facteur irreductible grand}. Let $(Q_i)_{i\geq 0}$ be the (infinite) sequence of factors $(\tP_j)_{j\geq i_0}$ reordered by increasing height, without repetition. By Lemma~\ref{lem: existence facteur irreductible grand}, we may assume $i_0$ large enough so that $\normH{Q_i} > 1$ for each~$i$, as well as $|Q_i(\xi)| \leq \normH{Q_i}^{-\homega}$. This sequence clearly satisfies the first assertions \eqref{enum: property (Q_i) item 1} to \eqref{enum: property (Q_i) item 6}, the third one coming from \eqref{eq lem: existence facteur irr grand eq 1} together with \eqref{eq: encadrement omega(P_i)}.

Now, let $i\geq 0$ and let $j\geq i_0$ be an index such that $Q_i$ divides $P_j$. Since we have $\normH{Q_i}\ll \normH{P_j}$ by Gelfond's lemma, the index $j$ tends to infinity as $i$ tends to infinity. Then, estimate \eqref{eq lem: existence facteur irr grand eq 2} can be rewritten as
\begin{equation}
\label{eq proof: inter 1 pour controler H(P_j)}
|P_j(\xi)|^{-1}\normH{P_j}^{-n-d+3} \ll |Q_i(\xi)|^{-1}\normH{Q_i}^{-n-d+3} = \normH{Q_i}^{\omega(Q_i)-n-d+3}.
\end{equation}
Using $|P_j(\xi)|^{-1} \gg \normH{P_j}^{\homega}$ and $\homega > 2n-d$, we get, for each large enough $i$,
\begin{equation*}
\normH{P_j}^{n-2d+3} \leq \normH{Q_i}^{\omega(Q_i)-n-d+3},
\end{equation*}
which is equivalent to \eqref{eq: premiere estimation Pj vs Qi}. So, assertion~\eqref{enum: H(P_j) controle par H(Q_i)} holds assuming $i_0$ large enough.

It remains to prove assertion~\eqref{enum: H(Q_i+1) controle par H(Q_i)}. Note that this is trivial if $\omega_n(\xi) = \infty$. Let us assume that $\omega_n(\xi) < \infty$ and fix a small $\ee > 0$ to be chosen later. For each pair $(i,j)$ as above with $j \geq i_0$ large enough as a function of $\ee$, we have $\omega(P_j) > \homega_n(\xi)-\ee/2$ and $\omega(Q_i) < \omega_n(\xi)+\ee/2$, and thus \eqref{eq proof: inter 1 pour controler H(P_j)} yields
\begin{equation*}
\normH{P_j}^{\homega_n(\xi)-\ee-n-d+3} \leq \normH{Q_i}^{\omega_n(\xi)+\ee-n-d+3},
\end{equation*}
for each $i \geq 0$ and each $j\geq i_0$ such that $Q_i$ divides $P_j$. We define $k$ as the largest index such that
\begin{equation*}
\normH{P_k} \leq \normH{Q_i}^{\theta(\ee)},\quad \textrm{where } \theta(\ee) = \frac{\omega_n(\xi)+\ee-n-d+3}{\homega_n(\xi)-\ee-n-d+3}.
\end{equation*}
Since $\normH{P_j} \leq \normH{Q_i}^{\theta(\ee)}$, by maximality of $k$ we have $i_0\leq j\leq k$. Let $\ell$ be such that $Q_\ell$ divides $P_{k+1}$. We find
\begin{equation*}
\normH{P_k} \leq \normH{Q_i}^{\theta(\ee)} < \normH{P_{k+1}} \leq \normH{Q_\ell}^{\theta(\ee)},
\end{equation*}
and therefore $\ell \geq i+1$. On the other hand, since by Gelfond's lemma we have $\normH{Q_\ell} \ll \normH{P_{k+1}}$, we deduce from \eqref{eq: H(P_i+1) controle par H(P_i)} that
\begin{equation*}
\normH{Q_{i+1}} \leq \normH{Q_\ell} \ll \normH{P_{k+1}} \ll \normH{P_k}^{\omega_n(\xi)/\homega}\leq \normH{Q_i}^{\omega_n(\xi)\theta(\ee) /\homega}.
\end{equation*}
We now choose $\ee>0$ small enough so that
\begin{equation*}
\theta(\ee) < \frac{\omega_n(\xi)-n-d+3}{\homega -n-d+3}.
\end{equation*}
This is possible since $\homega < \homega_n(\xi)$, and it yields \eqref{eq: estimation Q_i+1 par Q_i} for each $i\geq 0$, assuming that $i_0$ is large enough.
\end{proof}

\section{On the dimension of some polynomial subspaces}
\label{section: espaces V_N}

We start by introducing some families of vector spaces spanned by polynomials, and we study their dimensions.

\begin{Def}
\label{Def: fonction g_A(N)}
Let $k\geq n$ be an integer and let $\cA$ be a subset of $\bR[X]_{\leq n}$. We~define
\begin{align*}
\cB_k(\cA) &= \big\{ Q, XQ,\dots, X^{k-\deg(Q)}Q \,;\, Q\in\cA\setminus\{0\} \big\} \subset \bR[X]_{\leq k}, \\
V_k(\cA) &= \Vect[\bR]{\cB_k(\cA)}, \\
g_\cA(k) &= \dim V_k(\cA).
\end{align*}
\end{Def}

The spaces $V_k(\cA)$ play the role of the spaces $\cU^{k}(\cA)$ in \cite[\S3]{poels2022simultaneous} (for simultaneous approximation). We obtain analog properties. Note that if $\cA$ contains at least one non-zero polynomial, then
\begin{equation}
\label{eq: suite V_k(A) strict. croissante}
V_n(\cA) \varsubsetneq V_{n+1}(\cA) \varsubsetneq\cdots.
\end{equation}

The goal of this section is to prove the following result. We could not find a reference for the proposition below.

\begin{Prop}
\label{Cor: V_2n-k = tout l'espace}
Let $k$ be an integer with $0\leq k \leq n$, and let $\cA$ be a set of $k+1$ linearly polynomials of $\bR[X]_{\leq n}$. Suppose that the $\gcd$ of the elements of $\cA$ is $1$ (in~other words, the ideal spanned by $\cA$ is $\bR[X]$). Then
\begin{equation}
\label{Cor: V_N = tout l'espace}
V_{2n-k}(\cA) = \bR[X]_{\leq 2n-k}.
\end{equation}
\end{Prop}

The case $k=1$ is a classical result (it is implied by the fact that the resultant of two coprime polynomials is non-zero). The proof of Proposition~\ref{Cor: V_2n-k = tout l'espace} is given at the end of the section. Recall that a function $f : \{n,n+1, \ldots \} \to \bR$ is \emph{concave} if for any $i > n$, it satisfies
\[
f(i)-f(i-1) \geq f(i+1)-f(i).
\]
The next result is a dual version of \cite[Prop.\,3.1]{poels2022simultaneous} (where we deal with simultaneous approximation to the successive powers of $\xi$).

\begin{Lem}
Let $\cA \neq \{0\}$ be a non-empty subset of $\bR[X]_{\leq n}$. The function $g_\cA$ is concave and (strictly) increasing on $\{n,n+1,\ldots\}$.
\end{Lem}

\begin{proof}
The series of inclusions \eqref{eq: suite V_k(A) strict. croissante} shows that the function $g_\cA$ is increasing on $\{n,n+1,\dots\}$. For simplicity, we write $V_i=V_{i}(\cA)$ and $\cB_{i} = \cB_{i}(\cA)$ for each $i \geq n$. Given an integer $i\geq n$ we have $XV_{i}\subset V_{i+1}$, and we set
\[
h(i):= \dim \big(V_{i+1}/XV_{i}\big) = g_\cA(i+1)-g_\cA(i).
\]
We have to prove that $h$ is decreasing on $\{n,n+1,\cdots\}$. Fix $i\geq n+1$ and consider the linear map $\pi : V_i \to V_{i+1}/XV_{i}$ defined by $\pi(P) = P + XV_{i}$. Since $\cB_{i}\cup X\cB_{i} = \cB_{i+1}$, we have $V_i + XV_i = V_{i+1}$. So $\pi$ is surjective, and consequently $\textrm{Im } \pi = V_{i+1}/XV_{i}$ is isomorphic to $V_i / \ker \pi$.
On the other hand, $XV_{i-1}\subset V_{i} \cap XV_{i} \subset \ker\pi$, so $XV_{i-1}$ is subspace of $\ker \pi$. Hence
\begin{equation*}
h(i-1) = \dim \big(V_{i}/XV_{i-1}\big) \geq \dim\big(V_i / \ker \pi \big) = \dim \big(V_{i+1}/XV_{i}\big) = h(i).\qedhere
\end{equation*}
\end{proof}

\begin{Lem}
\label{lem : dim V_n+j(P,Q)}
Let $P,Q \in \bR[X]_{\leq n}$ be two coprime polynomials. Then, we have
\begin{equation*}
\dim V_{n+j}(P,Q) \geq 2(j+1),
\end{equation*}
for each $j\in\{0,\dots,n-1\}$. In particular $V_{2n-1}(P,Q) = \bR[X]_{\leq 2n-1}$.
\end{Lem}

\begin{proof}
Let $p$ (\resp $q$) denote the degree of $P$ (\resp of $Q$). There exist $\alpha,\beta\in\bR$ such that the polynomial $\tP := P(X)(X-\alpha)^{n-p}$ and $\tQ:= Q(X)(X-\beta)^{n-q}$ are coprime (and of degree exactly $n$). Fix $j\in\{0,\dots,n-1\}$. Since $\tP$ and $\tQ$ are coprime and $j<n$, the linear map
\begin{align*}
\bR[X]_{\leq j}\times \bR[X]_{\leq j} & \longrightarrow \bR[X]_{\leq n+j} \\
(R,S) & \longmapsto R\tP + S\tQ
\end{align*}
is injective, so its image $V_{n+j}(\tP,\tQ) \subset V_{n+j}(P,Q)$ has dimension $2(j+1)$.
\end{proof}

\begin{proof}[Proof of Proposition~\ref{Cor: V_2n-k = tout l'espace}]
For simplicity, we write $g=g_\cA$. Recall that $\cA$ has cardinality $k+1$, so that $g(n)\geq \textrm{card}(\cA) = k+1$. If $k=n$, then \eqref{Cor: V_N = tout l'espace} is automatic (since in that case $\cA$ contains a basis of $\bR[X]_{\leq n}$). So, we may assume that $k< n$. We first prove that for each sufficiently large $m$, we have
\begin{equation}
\label{eq : il existe m tel que V_m = tout}
V_{m}(\cA) = \bR[X]_{\leq m}.
\end{equation}
Indeed, since the ideal spanned by $\cA$ is $\bR[X]$, there exists an integer $\ell\geq n$ such that $1\in V_\ell(\cA)$. Let $P$ be a non-zero element in $\cA$ of degree $d$, and set $m=\ell+d$. Then $V_{m}(\cA)$ contains $\bR[X]_{\leq d}$, as well as the polynomials $P,XP,\ldots,X^\ell P$. We easily deduce \eqref{eq : il existe m tel que V_m = tout}.

By contradiction, suppose that \eqref{Cor: V_N = tout l'espace} does not hold, \ie
\begin{equation}
\label{eq proof:Cor: V_N = tout l'espace: by contradiction}
g(2n-k) \leq 2n-k.
\end{equation}
We distinguish two cases. Suppose first that $g(2n-k)-g(2n-k-1) \geq 2$. By concavity, then $g(j)-g(j-1)\geq 2$ for each $j$ with $n < j \leq 2n-k$, and we deduce that
\begin{equation*}
g(2n-k) \geq g(n)+2(n-k) \geq k+1 + 2(n-k) = 2n-k+1,
\end{equation*}
since $g(n) \geq \textrm{card}(\cA) = k+1$. This contradicts \eqref{eq proof:Cor: V_N = tout l'espace: by contradiction}, so $g(2n-k)-g(2n-k-1) \leq 1$. Since the function $g$ is increasing and concave, it is linear with slope $1$ on
\[
\{2n-k,2n-k+1,\dots\}.
\]
Choosing $m > 2n-k$ such that \eqref{eq : il existe m tel que V_m = tout} holds, we obtain by \eqref{eq proof:Cor: V_N = tout l'espace: by contradiction}
\begin{equation*}
m+1 = g(m) = g(2n-k)+ m - (2n-k) \leq m,
\end{equation*}
a contradiction. Hence \eqref{Cor: V_N = tout l'espace} holds.
\end{proof}

\begin{figure}[htb]
\begin{tikzpicture}[xscale=0.5,yscale=0.4, line cap=round,line join=round,>=triangle 45,x=1cm,y=1cm]
\clip(-5,-2) rectangle (20,15);
\draw[-stealth, semithick] (-0.15,0)--(18,0) node[below]{$i$};
\draw[-stealth, semithick] (0,-0.15)--(0,14.5) node[left]{$g(i)$};

\draw [thick] (0,1)-- (1,5);
\draw [thick] (1,5)-- (2.66,8.059);
\draw [thick] (2.66,8.059)-- (7,10.5);
\draw [thick] (7,10.5)-- (14,12);
\draw [dashed] (7,13.34)-- (7,-0.5);

\node[left] at (-0,1.7) {$\geq \textrm{card}(A)$};
\node[above] at (4,12) {slope $\geq 2$};
\node[above] at (10,12) {slope $= 1$};
\node[below] at (7,0) {$2n-\ell$};
\node[below] at (14,0) {$2n-1$};
\node[left] at (0,13) {$2n$};

\draw [dash pattern=on 2pt off 3pt](0,12)-- (14,12);
\draw [dash pattern=on 2pt off 3pt] (14,12)-- (14,0);

\node[below] at (0,-0.15) {$i=n$};
\node[draw,circle,inner sep=1.25pt,fill] at (0,1) {};
\node[draw,circle,inner sep=1.25pt,fill] at (1,5) {};
\node[draw,circle,inner sep=1.25pt,fill] at (2.666,8.059) {};
\node[draw,circle,inner sep=1.25pt,fill] at (7,10.5) {};
\node[draw,circle,inner sep=1.25pt,fill] at (14,12) {};
\node[draw,circle,inner sep=1.25pt,fill] at (0,12) {};
\node[draw,circle,inner sep=1.25pt,fill] at (14,0) {};
\node[draw,circle,inner sep=1.25pt,fill] at (7,0) {};
\node[draw,circle,inner sep=1.25pt,fill] at (0,0) {};
\end{tikzpicture}
\caption{Graph of the piecewise linear function interpolating the
values $g(i)=\dim V_i(\cA)$ at integers $i\in\{n,\dots,2n-1\}$.}
\label{fig:1}
\end{figure}

\section{Proof of Theorem~\ref{Thm : main d=2} (case \texorpdfstring{$d=2$}{d2})}
\label{Section: cas d = 2}

In this section, we deal with the case $d=2$ to prove Theorem~\ref{Thm : main d=2}, namely that $\homega_3(\xi)\leq 2+\sqrt 5 = 4.23\cdots$ and $\homega_n(\xi) \leq 2n-2$ for each $n\geq 4$. The estimate $\homega_n(\xi) \leq 2n-2$ was already known for $n\geq 10$, however for $n=4,\dots,9$ it is a new result. For $n=3$, our bound improves on the bound $\homega_3(\xi)\leq 3+\sqrt 2 = 4.41\cdots$ due to Bugeaud and Schleischitz \cite{bugeaudSchleischitz2016Uniform}. Moreover, our proof does not require Marnat-Moshchevitin's inequality \cite{marnat2018optimal}.

\begin{proof}[Proof of Theorem~\ref{Thm : main d=2}]
Suppose that $\homega_n(\xi) > 2n-2$, and fix a real number $\homega$ such that
\[
\homega_n(\xi) > \homega > 2n-2.
\]
Let $(P_i)_{i\geq 0}$ be a sequence of minimal polynomials associated to $n$ and $\xi$ as in Section~\ref{section: minimal pol}. According to Lemma~\ref{lem: existence facteur irreductible grand} (with $d=2$) there exists an index $i_0\geq 0$ such that $P_i$ has degree $n$ and is irreducible for each $i\geq i_0$. Consequently, up to a finite number of terms, the sequence $(P_i)_{i\geq 0}$ coincides with the sequence $(Q_i)_{i\geq 0}$ of Proposition~\ref{prop: existence Qi}. Let $I$ denotes the set of indices $i\geq i_0+1$ such that $P_{i-1}$, $P_i$ and $P_{i+1}$ are linearly independent. By Lemmas~\ref{Lem: P_iH_i+1 = P_j-1H_j} and~\ref{lem: I infini}, the set $I$ is infinite, and for any consecutive $i<j$ in $I$, we have
\[
\normH{P_{i+1}}|P_i(\xi)| \asymp \normH{P_j}|P_{j-1}(\xi)|.
\]
Furthermore, the irreducible polynomials $P_i$ and $P_{i+1}$ are also coprime since $\normH{P_i} < \normH{P_{i+1}}$ and $|P_i(\xi)| > |P_{i+1}(\xi)|$. Lemma~\ref{lem: estimation classique du resultant} yields
\begin{align*}
1 \leq |\Res(P_i,P_{i+1})| \ll \normH{P_i}^{n-1}\normH{P_{i+1}}^n|P_{i}(\xi)| &\ll \normH{P_i}^{n-1}\normH{P_{j}}^n|P_{j-1}(\xi)| \\
& \ll \normH{P_i}^{n-1}\normH{P_{j}}^{n-\homega}.
\end{align*}
We deduce that
\begin{equation}
\label{eq proof: estimation theta}
\normH{P_{j}} \leq \normH{P_i}^\theta \quad \textrm{where } \theta = \frac{n-1}{\homega-n}.
\end{equation}
Let $h<i<j$ be consecutive indices in $I$. We have the following configuration
\[
\Vect[\bR]{P_h,P_{h+1}} = \Vect[\bR]{P_{i-1},P_{i}} \neq \Vect[\bR]{P_{i},P_{i+1}},
\]
so $P_{h}, P_{h+1}, P_{i+1}$ are linearly independent. Proposition~\ref{Cor: V_2n-k = tout l'espace} combined with Lemma~\ref{lem : dim V_n+j(P,Q)} implies that
\begin{equation*}
\Big(\bR[X]_{\leq n-2}P_{h}\oplus \bR[X]_{\leq n-2}P_{h+1}\Big) + \bR[X]_{\leq n-2}P_{i+1} = \bR[X]_{\leq 2n-2}.
\end{equation*}
Choose $k\in\{0,\dots,n-2\}$ such that $\big(P_{h},\dots,X^{n-2}P_{h},P_{h+1},\dots,X^{n-2}P_{h+1}, X^kP_{i+1}\big)$ is a basis of $\bR[X]_{\leq 2n-2}$. We denote by $M$ the matrix of this basis expressed in the canonical basis $(1,X,\dots,X^{2n-2})$. Estimating $\det(M)$ as in the proof of Lemma~\ref{lem: estimation classique du resultant} (in other words, for $\ell=2,\dots,2n-2$, we add to the first row of $M$ the $\ell$-th row multiplied by $\xi^{\ell-1}$), we get the estimates
\begin{equation*}
1 \leq \big|\det\big(M)\big| \ll |P_{h}(\xi)|\normH{P_{h}}^{n-2}\normH{P_{h+1}}^{n-1}\normH{P_{i+1}}.
\end{equation*}
Now, since $\normH{P_{h+1}}^{n-1}|P_{h}(\xi)| \asymp \normH{P_{h+1}}^{n-2}\normH{P_i}|P_{i-1}(\xi)| \ll \normH{P_i}^{n-1-\homega}$, we deduce that
\begin{equation}
\label{eq proof: estimation tau}
\normH{P_i}^{\homega-n+1} \ll \normH{P_{h}}^{n-2}\normH{P_{j}}.
\end{equation}
For consecutive $i<j$ in $I$, define $\tau_i\in(0,1)$ by
\[
\normH{P_i} = \normH{P_j}^{\tau_i}
\]
and set $\tau = \limsup_{i\in I, i\to\infty} \tau_i \in [0,1]$. Let $h<i<j$ be consecutive indices in $I$ as previously. By \eqref{eq proof: estimation tau}, we obtain
\[
\homega-n+1 \leq (n-2)\tau_h + \frac{1}{\tau_i} +o(1) \leq (n-2)\tau + \frac{1}{\tau_i} +o(1).
\]
We infer that
\begin{equation}
\label{eq proof: majoration homega_3: eq 2}
p(\tau) \geq 0, \quad \textrm{where } p(t) = (n-2)t^2-(\homega-n+1)t+1.
\end{equation}
Note that
\begin{equation*}
p(0) = 1, \qquad p\Bigl(\frac{1}{n-2}\Bigr) = \frac{2n-2-\homega}{n-2} < 0 \AND p(1) = 2n-2-\homega < 0.
\end{equation*}
We deduce that $p$ has one root $\alpha\in(0,1/(n-2))$ and one root larger than $1$. Since $\tau\in[0,1]$ and $p(\tau)\geq 0$, we obtain $\tau \leq \alpha$. Combined with the estimate $\normH{P_i} = \normH{P_j}^{\tau_i} \ll \normH{P_i}^{\theta\tau_i}$ valid
for any $i\in I$ (this is a consequence of \eqref{eq proof: estimation theta}), this leads to
\begin{equation}
\label{eq proof: majoration homega_3: eq 3}
1 \leq \theta\tau \leq \theta\alpha < \frac{n-1}{(n-2)^2}.
\end{equation}
We easily check that this is impossible when $n\geq 4$ (the right-hand side is strictly less than $1$), thus $\homega_n(\xi)\leq 2n-2$ for each $n\geq 4$.

We now deal with the case $n=3$. Suppose by contradiction that $\homega_3(\xi) > 2+\sqrt 5$ and choose $\homega$ such that
\[
\homega_3(\xi) > \homega > 2+\sqrt 5.
\]
The polynomial $p$ from \eqref{eq proof: majoration homega_3: eq 2} becomes $p(t) = t^2-(\homega-2)t+1$. Denote by $\alpha$ its smallest root, and by $\beta = (\sqrt 5 -1)/2$ the smallest root of the polynomial $t^2-\sqrt 5 t + 1$. We find
\begin{equation*}
0 = \beta^2-\sqrt 5 \beta + 1 > \beta^2-(\homega-2 )\beta + 1 = p(\beta),
\end{equation*}
hence $\alpha < \beta$. Combined with $\theta = 2/(\homega-3) < 1/\beta$, this implies that $\theta \alpha < 1$, which contradicts \eqref{eq proof: majoration homega_3: eq 3}. It follows that $\homega_3(\xi) \leq 2+\sqrt 5$.
\end{proof}


\begin{thebibliography}{99}
\bibitem{arbourRoyCriterionDegreTwo} B. Arbour \& D. Roy – “A Gel’fond type criterion in degree two”, Acta Arith. 111 (2004), no. 1, p. 97–103. \url{https://doi.org/10.4064/aa111-1-7}
\bibitem{bourbakiAlgebre} N. Bourbaki – Algebra I, Springer-Verlag, New York, 1989. 
\bibitem{brownawell1974sequences} W. D. Brownawell – “Sequences of Diophantine approximations”, J. Number Theory 6 (1974), no. 1, p. 11–21. \url{https://doi.org/10.1016/0022-314X(74)90004-3}
\bibitem{bugeaud2004approximation} Y. Bugeaud – Approximation by algebraic numbers, Cambridge Tracts in Math., vol. 160, Cambridge University Press, Cambridge, 2004. \url{https://doi.org/10.1017/CBO9780511542886}
\bibitem{bugeaud2015exponents} \textemdash{}, “Exponents of Diophantine approximation”, in Dynamics and analytic number theory (D. Badziahin, A. Gorodnik \& N. Peyerimhoff, eds.), London Math. Soc. Lect. Note Ser., vol. 437, Cambridge University Press, Cambridge, 2016, p. 96–135. \url{https://doi.org/10.1017/9781316402696.003}
\bibitem{bugeaud2005exponentsSturmian} Y. Bugeaud \& M. Laurent – “Exponents of Diophantine approximation and Sturmian continued fractions”, Ann. Inst. Fourier (Grenoble) 55 (2005), no. 3, p. 773–804. \url{https://doi.org/10.5802/aif.2114}
\bibitem{bugeaud2010transfer} \textemdash{}, “On transfer inequalities in Diophantine approximation, II”, Math. Z. 265 (2010), no. 2, p. 249–262. \url{https://doi.org/10.1007/s00209-009-0512-0}
\bibitem{bugeaudSchleischitz2016Uniform} Y. Bugeaud \& J. Schleischitz – “On uniform approximation to real numbers”, Acta Arith. 175 (2016), no. 3, p. 255–268.  
\bibitem{davenport1967approximation} H. Davenport \& W. Schmidt – “Approximation to real numbers by quadratic irrationals”, Acta Arith. 13 (1967), no. 2, p. 169–176. \url{https://doi.org/10.4064/aa-13-2-169-176}
\bibitem{davenport1969approximation} H. Davenport \& W. Schmidt – “Approximation to real numbers by algebraic integers”, Acta Arith. 15 (1969), no. 4, p. 393–416. \url{https://doi.org/10.4064/aa-15-4-393-416}
\bibitem{Hodge} W. Hodge \& D. Pedoe – Methods of algebraic geometry, Cambridge University Press, Cambridge, 1947. 
\bibitem{laurent2003simultaneous} M. Laurent – “Simultaneous rational approximation to the successive powers of a real number”, Indag. Math. (N.S.) 14 (2003), no. 1, p. 45–53. \url{https://doi.org/10.1016/S0019-3577(03)90070-X}
\bibitem{marnat2018optimal} A. Marnat \& N. Moshchevitin – “An optimal bound for the ratio between ordinary and uniform exponents of Diophantine approximation”, Mathematika 66 (2020), no. 3, p. 818–854. \url{https://doi.org/10.1112/mtk.12045}
\bibitem{moshchevitin2007best} N. G. Moshchevitin – “Best Diophantine approximations: the phenomenon of degenerate dimension”, in Surveys in geometry and number theory: reports on contemporary Russian mathematics, London Math. Soc. Lecture Note Ser., vol. 338, Cambridge University Press, Cambridge, 2007, p. 158–182. \url{https://doi.org/10.1017/CBO9780511721472.006}
\bibitem{poels2017exponents} A. Poëls – “Exponents of Diophantine approximation in dimension 2 for numbers of Sturmian type”, Math. Z. 294 (2020), no. 3, p. 951–993. \url{https://doi.org/10.1007/s00209-019-02280-2}
\bibitem{poelsExpoGeneralClass2021} \textemdash{}, “Exponents of Diophantine approximation in dimension two for a general class of numbers”, Mosc. J. Comb. Number Theory 11 (2022), no. 1, p. 37–69. \url{https://doi.org/10.2140/moscow.2022.11.37}
\bibitem{poels2022simultaneous} A. Poëls \& D. Roy – “Simultaneous rational approximation to successive powers of a real number”, Trans. Amer. Math. Soc. 375 (2022), no. 09, p. 6385–6415. 
\bibitem{poelsroy2022parametric} \textemdash{}, “Parametric geometry of numbers over a number field and extension of scalars”, Bull. Soc. math. France 151 (2023), p. 257–303. 
\bibitem{PhDMartin2019} M. Rivard-Cooke – “Parametric geometry of numbers”, PhD Thesis, University of Ottawa, 2019, https://ruor.uottawa.ca/handle/10393/38871. 
\bibitem{roy2003approximation} D. Roy – “Approximation simultanée d’un nombre et de son carré”, Comptes Rendus Mathématique 336 (2003), no. 1, p. 1–6. 
\bibitem{roy2004approximation} \textemdash{}, “Approximation to real numbers by cubic algebraic integers I”, Proc. London Math. Soc. (3) 88 (2004), no. 1, p. 42–62. 
\bibitem{roy2007two} \textemdash{}, “On two exponents of approximation related to a real number and its square”, Canad. J. Math. 59 (2007), no. 1, p. 211–224. 
\bibitem{schleischitz2017some} J. Schleischitz – “Some notes on the regular graph defined by Schmidt and Summerer and uniform approximation”, JP J. Algebra Number Theory Appl. 39 (2017), no. 2, p. 115–150. \url{https://doi.org/10.17654/NT039020115}
\bibitem{schleischitz2017uniformPoly} \textemdash{}, “Uniform Diophantine approximation and best approximation polynomials”, Acta Arith. 185 (2018), no. 3, p. 249–274. \url{https://doi.org/10.4064/aa170901-4-7}
\bibitem{schleischitz2024uniform} \textemdash{}, “Uniform dual approximation to Veronese curves in small dimension”, 2024, arXiv: 2405.10086. \url{https://arxiv.org/abs/2405.10086}
\bibitem{schmidt1996diophantine} W. M. Schmidt – Diophantine approximation, Lect. Notes in Math., vol. 785, Springer-Verlag, Berlin, 1980. 
\bibitem{Schmidt} \textemdash{}, Diophantine approximations and Diophantine equations, Lect. Notes in Math., vol. 1467, Springer-Verlag, Berlin, 1991. \url{https://doi.org/10.1007/BFb0098246}
\end{thebibliography}
\end{document}
